\chapter{Summary and Future Work}\label{ch:summary}

This report addressed the gap between geohazard uncertainty quantification and computational energy profiling. By coupling the SynxFlow hydrodynamic solver with the Alumet telemetry framework, this work transitioned from relying on execution time as a proxy to directly measuring process-isolated energy consumption for stochastic flood simulations. This chapter summarizes the key findings, acknowledges experimental limitations, and outlines future research directions.

%%%%%%%%%%%%%%%%%%%%%%%%%%%%%%%%%%%%%%%%%%%%%%%%%%%%%%%%%%%%%%%%%%%%%%%%
\section{Summary}\label{sec:summary}

This research investigated how input data uncertainty and execution environment volatility impact the energy predictability of GPU-accelerated flood solvers. The core research questions from Section~\ref{sec:research-questions} were addressed as follows:
  
\begin{itemize}
    \item[\textit{RQ1}] \textbf{Methodologies for process-level energy attribution:} A custom Python framework was developed to integrate SynxFlow with Alumet. By utilizing dynamic Process ID tracking and scaling hardware power by utilization percentages, this methodology isolated the solver's true energy footprint from shared cluster noise.
    
    \item[\textit{RQ2}] \textbf{Execution time as an energy proxy:} Preliminary results indicated that execution time may be an unreliable standalone proxy for computational energy in this context. Welch's t-tests suggested that altering the computational grid's spatial resolution can shift the hardware's dynamic power intensity. Time-to-solution might not fully capture this energy variability, as power rates appeared to fluctuate based on grid-to-GPU alignment.

    \item[\textit{RQ3}] \textbf{Uncertainty propagation to energy variability:} Initial observations suggested that as spatial input uncertainty ($\sigma$) increases, both the average computational energy and run-to-run variance show a tendency to expand. Rougher terrain may trigger adaptive timestepping, potentially producing right-skewed energy distributions. This indicates that highly uncertain DEMs could lead to highly uncertain energy footprints, suggesting that single deterministic profiles might underestimate the energy requirements of stochastic ensembles.
\end{itemize}
%%%%%%%%%%%%%%%%%%%%%%%%%%%%%%%%%%%%%%%%%%%%%%%%%%%%%%%%%%%%%%%%%%%%%%%%

\section{Limitations}
\label{sec:limitations}

While establishing a foundational link between input uncertainty and energy consumption, this study has several boundaries. First, the simplified error modeling relies on a zero-mean Gaussian noise model ($\epsilon \sim \mathcal{N}(0, \sigma^2)$), which assumes isotropic errors and oversimplifies the spatially correlated and biased errors of real-world remote sensing. Second, the scope of the scenario is limited, as the findings stem from a single flood simulation. Consequently, the observed non-linear hydrodynamic effects may not uniformly apply to varying catchments or urban floodplains.

Additionally, due to computational expense, the ensemble sizes were restricted to 100 iterations. While this sample size captures primary trends, it may be insufficient to fully stabilize extreme tail behaviors within the energy distributions. Finally, the observed power shifts are inherently tied to the hardware specificity of the Gaia compute cluster and its NVIDIA GPUs, meaning these exact energy dynamics may vary across different computing architectures.

%%%%%%%%%%%%%%%%%%%%%%%%%%%%%%%%%%%%%%%%%%%%%%%%%%%%%%%%%%%%%%%%%%%%%%%%
\section{Future Work}
\label{sec:future-work}

The findings of this thesis open several promising directions for further research. First, future studies must improve the error model by incorporating spatially correlated noise alongside a negative bias shift. Because spaceborne sensors frequently overestimate ground heights by capturing vegetation and buildings instead of bare earth \cite{Meadows2024}, an accurate error model must simulate the removal of this positive sensor bias by lowering the elevation values in obstructed areas.

Furthermore, the topographic scope of the framework should be expanded. Testing the methodology on diverse flood cases and different DEMs will help determine if specific terrain features inherently demand higher adaptive energy costs. Finally, scaling the Monte Carlo simulations to larger ensembles is necessary to observe if the variability in energy consumption fully stabilizes. Capturing this stabilization is critical for establishing reliable, carbon-aware HPC scheduling limits.